\documentclass[11pt]{article}
\usepackage[T1]{fontenc}
\usepackage[utf8]{inputenc}
\usepackage{lmodern}
\usepackage[margin=1in]{geometry}
\usepackage{amsmath,amssymb,mathtools}
\usepackage{booktabs,longtable,array}
\usepackage{fancyvrb}
\usepackage[hidelinks]{hyperref}
\usepackage{microtype}
\usepackage{graphicx}
\usepackage{adjustbox}
\pdfinfoomitdate=1
\pdftrailerid{}
\pdfsuppressptexinfo=15
\setlength{\parindent}{0pt}
\setlength{\parskip}{0.6em}
\setlength{\emergencystretch}{3em}
\renewcommand{\arraystretch}{1.25}
\setlength{\tabcolsep}{4pt}
\sloppy
\graphicspath{{Revision/universe_model/}}
\title{UM-1: a 4+4 model of the Universe with Einstein-Lovelock gravity, the dirac16complex field and an x0-generated mass spectrum}
\author{dirac16complex project (with Claude Opus 5.5)}
\date{}
\begin{document}
\maketitle
\tableofcontents
\newpage

Version 1.2 of 2026-10-17: version 1 plus the two changes that the solution forces (marked CHANGED in sections 4.1 and 8; the reasons are in \texttt{SOLUTION.\allowbreak{}md} section 3), and (1.2) the corrections of the analysis \texttt{ANALYSIS.\allowbreak{}md} to sections 1, 4.4, 6 and 8 and the new section 9; the fields, the action and the parameters are those of version 1.1. Written at the author's order of 2026-10-17 ("create a new model of the Universe"). This document DEFINES the model. Its field equations are derived in \texttt{FIELD\_\allowbreak{}EQUATIONS.\allowbreak{}md} (scripts in \texttt{field\_\allowbreak{}equations/\allowbreak{}}), an approximate solution is computed in \texttt{SOLUTION.\allowbreak{}md} (scripts in \texttt{solution/\allowbreak{}}), and the origin of every ingredient is in \texttt{PROVENANCE.\allowbreak{}md}. Every statement below is labelled as one of: DERIVED (computed in this repository, with the file named), TAKEN (from the Revision record, file named), STANDARD (textbook physics, used as input), ASSUMED (a choice of this model), INVENTED (a new phenomenological ingredient of this model) or OPEN.

\section*{1. Summary}
\addcontentsline{toc}{section}{1. Summary}

The model lives on an eight-dimensional space with the author's coordinates $x_0, x_1, \dots, x_7$: $x_0$ is the hidden space-like direction, $x_1, x_2, x_3$ are ordinary space, $x_4$ is the time and $x_5, x_6, x_7$ are three extra time-like directions that deflate exponentially. Gravity is Einstein-Lovelock gravity in eight dimensions with the author's couplings. The dirac16complex field $\Psi$, a 16-component spinor of $Cl(4,4)$ coupled to gravity through the canonical spin connection, is the dark-matter candidate: its heavy odd quanta conditionally, its massless fold band as dark radiation, and none of its states as dark energy (\texttt{ANALYSIS.\allowbreak{}md}, section 9 below). Photons, neutrinos, electrons, up and down quarks, the gluons (required for consistency with QCD), and the $W^\pm$ and $Z$ bosons are eight-dimensional fields of the gauge group $SU(3)_c \times SU(2)_L \times U(1)_Y$. There is NO Higgs field and no Higgs particle.

The central new idea (the author's): every field is a function of $x_0$ that is periodic with the period $\ell_0 = \pi/(6H)$ of the metric. Expanding a field in the eigenfunctions of its $x_0$ operator turns the eight-dimensional wave equation into a tower of four-dimensional fields whose masses are the eigenvalues. A mode that does not depend on $x_0$ is massless (graviton, photon, gluon). A mode localised in $x_0$ carries a hidden momentum $\Delta p_{x_0} \geq \hbar/(2\,\Delta x_0)$ and therefore an effective mass. The model makes this precise: the four-dimensional mass of every massive particle is the lowest eigenvalue of its $x_0$ operator on the periodic cell, with the boundary conditions at the two special points of the cell (the fold $z = \pi/2$ and the tip $z = 0$). CORRECTED in version 1.2 (\texttt{ANALYSIS.\allowbreak{}md}): the uncertainty relation bounds the hidden kinetic part of every such eigenvalue, and it bounds the mass itself for the $x_0$ excitations (every mode heavier than the gap $3H/(2f)$, including the dark matter). It does not bound the masses of the protected zero modes, which are localised at the fold yet massless, and the light masses ($W$, $Z$, $e$, $\nu$, $u$, $d$) come from the tip, not from localisation (section 9).

\section*{2. Space, coordinates and the periodic cell}
\addcontentsline{toc}{section}{2. Space, coordinates and the periodic cell}

\begin{itemize}
\item Coordinates (TAKEN, \texttt{Revision/\allowbreak{}SPEC.\allowbreak{}md} section 1): $z = 6Hx_0$ takes every real value; the metric depends on $x_0$ only through $\cot^2 z$ and $(\sin^2 z)^{1/6}$, both periodic with period $\pi$ in $z$, i.e. with period $\ell_0 = \pi/(6H)$ in $x_0$.
\item The cell (TAKEN, \texttt{Revision/\allowbreak{}kohn\_\allowbreak{}sham/\allowbreak{}ks-{}\allowbreak{}theory.\allowbreak{}json}, boundaryConditions): one period $0 < z < \pi$ consists of two halves, $0 < z < \pi/2$ and $\pi/2 < z < \pi$, joined at the FOLD $z = \pi/2$ ($\cos z = 0$, where $g_{00}$ and $\sqrt{|g|}$ vanish) and ending at the TIP $z \to 0, \pi$ ($\sin z \to 0$). The hidden coordinate $y = \ln|\sin z|/(6H) \in (-\infty, 0]$ measures the distance from the fold; $dy = \cot z\, dx_0$.
\item A field periodic in $x_0$ with period $\ell_0$ is a field on the cell. Its two halves are related by the fold condition $(1 - Q(\theta))\chi(0) = 0$, $Q(\theta) = \cos\theta\,\sigma_3 + \sin\theta\,\sigma_2$, with $\theta = 0$ (even) or $\theta = \pi$ (odd) (DERIVED in the Kohn-Sham record up to $\theta$). Periodicity in $x_0$ requires that the two tips of a cell ($z \to 0^+$ and $z \to \pi^-$) be identified; the proper distance to the tip is infinite, so the model replaces the tip by a regular cutoff at $y = -L$ (ASSUMED, as in the Kohn-Sham record; \texttt{tip\_\allowbreak{}convergence/\allowbreak{}} measures the dependence on $L$).
\end{itemize}

\section*{3. The metric as a power series}
\addcontentsline{toc}{section}{3. The metric as a power series}

The metric is written as a series of terms of decreasing magnitude (ASSUMED structure):

\[g_{MN} = \bar g_{MN} + \epsilon\, h^{(1)}_{MN} + \epsilon^2 h^{(2)}_{MN} + \dots\]

\begin{itemize}
\item Order 0 (the background, TAKEN from \texttt{Revision/\allowbreak{}SPEC.\allowbreak{}md} section 1, the corrected primordial metric, version 2):
\end{itemize}
\[\adjustbox{max width=\linewidth}{$\displaystyle \bar g = \mathrm{diag}\big(e^{-2a_4}b_4\cot^2 z,\ e^{2a_4},\ e^{2a_4},\ e^{2a_4},\ -1,\ -e^{-2a_4}b_4(\sin^2 z)^{1/6},\ -e^{-2a_4}b_4(\sin^2 z)^{1/6},\ -e^{-2a_4}b_4(\sin^2 z)^{1/6}\big)$}\]
with two GENERAL functions of the time, $a_4(x_4)$ (the observer's scale factor $a = e^{a_4}$) and $b_4(x_4) > 0$. The block $(x_0, x_5, x_6, x_7)$ has the common scale factor $f = e^{-a_4}b_4^{1/2}$ and the rate $h_4 = \dot f/f = u/2 - \dot a_4$, $u = \dot b_4/b_4$; the extra times deflate when $h_4 < 0$. The primordial history of the record is $a_4 = Hx_4$, $b_4 = e^{(4/3)Hx_4}$ ($h_4 = -H/3$). The model keeps $a_4$ and $b_4$ general and determines them from the field equations with all sources (section 7).
\begin{itemize}
\item Order 1: the cosmological perturbations, in the 3+1 decomposition of $(x_1, \dots, x_4)$: the scalar potentials $\Phi, \Psi_N$ (Newtonian gauge, $\delta g_{44} = -2\Phi$, $\delta g_{ij} = -2\Psi_N e^{2a_4}\delta_{ij}$), the tensor modes $h^{TT}_{ij}$, and the block perturbation $\delta b_4/b_4 = 2\sigma$ (a radion), all functions of $(x_1, \dots, x_4)$ times the flat $x_0$ zero-mode profile; the $x_0$-dependent parts of the sources drive $x_0$-dependent corrections, which are expanded in the same $x_0$ eigenfunctions as the fields and are heavy (section 5). The bookkeeping parameter $\epsilon$ is of the size of the CMB temperature anisotropy, $\epsilon \sim 10^{-5}$ (STANDARD).
\item Order 2 and beyond: the products of order-1 quantities and the $x_0$-heavy modes; not computed in version 1 (OPEN).
\end{itemize}

\section*{4. Fields and action}
\addcontentsline{toc}{section}{4. Fields and action}

\[\adjustbox{max width=\linewidth}{$\displaystyle S = S_{EL}[g] + S_\Psi[g, \Psi] + S_{SM}[g, A, \psi] + S_{tip}[g, A, \psi, \Sigma] + S_{portal}[\Psi, A, \psi]$}\]

\subsection*{4.1 Einstein-Lovelock gravity (TAKEN)}
\addcontentsline{toc}{subsection}{4.1 Einstein-Lovelock gravity (TAKEN)}

\[\adjustbox{max width=\linewidth}{$\displaystyle S_{EL} = \frac{1}{2\kappa}\int d^8x \sqrt{|g|}\,\Big(\sum_{k=1}^{3}\alpha_k\,\mathcal{L}_{(k)} - 2\Lambda\Big),\qquad \sum_k \alpha_k E_{(k)}{}^M{}_N + \Lambda\,\delta^M_N = \kappa\,T^M{}_N$}\]

with $\mathcal{L}_{(1)} = R$, $\mathcal{L}_{(2)}$ the Gauss-Bonnet and $\mathcal{L}_{(3)}$ the cubic Lovelock density, the author's weights $\alpha_k = w_k H^{2-2k}$, $w = (1, 1/32, 1/1152)$, and $\Lambda = -10H^2$ in the convention of \texttt{Revision/\allowbreak{}field\_\allowbreak{}equations\_\allowbreak{}a4/\allowbreak{}a4-{}\allowbreak{}equations.\allowbreak{}json}. The graviton is the $x_0$ zero mode of $g_{MN}$ and is massless.

CHANGED in version 1.1 (\texttt{SOLUTION.\allowbreak{}md} section 3, DERIVED there): with the author's weights the only static vacuum, $k_0 = 1/(4w_2) = 8$ with $\Lambda_0 = -24H^2$, has the effective-Newton factor $G(k_0) = -5 + 9w_3/(2w_2^2) = -1$, i.e. repulsive four-dimensional gravity, and $\Lambda = -10H^2$ admits no static block with a small cosmological term. Version 1.1 uses $w = (1, 1/32, 1/768)$ ($G(k_0) = 1$, the Einstein value) and $\Lambda = -24H^2 + \delta\Lambda$, and ASSUMES a block stabiliser (an extra potential for $k = e^{2a_4}/b_4$ with curvature above $k_0G(k_0)H^2/4 = 2H^2$, since the block mode of the homogeneous reduction has $\lambda^2 = k_0G(k_0)H^2/4$: attractive gravity makes the static block unstable). The origin of the stabiliser (the Casimir energy of the cell) is OPEN.

\subsection*{4.2 The dirac16complex field (TAKEN)}
\addcontentsline{toc}{subsection}{4.2 The dirac16complex field (TAKEN)}

\[\adjustbox{max width=\linewidth}{$\displaystyle S_\Psi = \int d^8x\sqrt{|g|}\,\Big[\tfrac12\big(\bar\Psi\gamma^M D_M\Psi - (D_M\bar\Psi)\gamma^M\Psi\big) - m\,\bar\Psi\Psi - \tfrac{\lambda}{2}(\bar\Psi\Psi)^2\Big],\quad D_M = \partial_M + \Omega_M,\ \Omega_M = \tfrac12\omega_M{}^{ab}S_{ab}$}\]

with $\bar\Psi = \Psi^\dagger C$, the author's $\gamma$ matrices (\texttt{Revision/\allowbreak{}algebra/\allowbreak{}gammas.\allowbreak{}json}) and the CANONICAL spin connection of the diagonal vielbein of $\bar g$, which is non-trivial: $\gamma^M\Omega_M = \tfrac32 H e^{a_4}b_4^{-1/2}\gamma^{x_0} + (u - \tfrac12\dot a_4)\gamma^{x_4}$ (DERIVED, \texttt{Revision/\allowbreak{}lead\_\allowbreak{}checks/\allowbreak{}}). $\Psi$ is a singlet of $SU(3)_c \times SU(2)_L \times U(1)_Y$ (ASSUMED: dark). Its quanta are of three kinds (DERIVED in the Kohn-Sham record): the fold-band quanta, $\varepsilon = e^{-a_4}|k|$ exactly, MASSLESS in four dimensions (the fold zero modes of the $x_0$ operator); the bulk-band quanta with $\varepsilon \geq m$ (massive; the $x_0$ momentum adds $(l\pi/(fL))^2$ to $m^2$); and the homogeneous condensate $S = \bar\Psi\Psi = S_0 e^{a_4}b_4^{-2}$.

\subsection*{4.3 The Standard-Model sector in eight dimensions (ASSUMED structure, STANDARD content)}
\addcontentsline{toc}{subsection}{4.3 The Standard-Model sector in eight dimensions (ASSUMED structure, STANDARD content)}

\begin{itemize}
\item Gauge fields $G^\alpha_M$ ($SU(3)_c$), $W^a_M$ ($SU(2)_L$), $B_M$ ($U(1)_Y$) with $S_{gauge} = -\frac14\int d^8x\sqrt{|g|}\, g^{MP}g^{NQ}\big(\frac{1}{\hat g_s^2}G^\alpha_{MN}G^\alpha_{PQ} + \frac{1}{\hat g^2}W^a_{MN}W^a_{PQ} + \frac{1}{\hat g'^2}B_{MN}B_{PQ}\big)$; the four-dimensional couplings are $g^2 = \hat g^2/V_\perp$ etc. with $V_\perp$ the proper volume of the cell and the extra-time block (section 5).
\item One generation of fermions, each an eight-dimensional spinor of $Cl(4,4)$ (16 components, the same algebra as $\Psi$): $L = (\nu, e)_L$ (doublet, $Y = -\frac12$), $e_R$ ($Y = -1$), $\nu_R$ ($Y = 0$), $Q = (u, d)_L$ (colour triplet, doublet, $Y = \frac16$), $u_R$ ($Y = \frac23$), $d_R$ ($Y = -\frac13$), with $S_\psi = \int d^8x\sqrt{|g|}\,\big[\bar\psi\gamma^M D_M\psi - \mu_\psi\,\bar\psi\psi\big]$, $D_M$ containing the same canonical spin connection and the gauge fields. Each species has a bulk mass $\mu_\psi$ (constant on the cell) and its own fold condition $\theta_\psi$ (INVENTED assignment): exactly as for $\Psi$, the reduced $x_0$ operator is the block $h = -(i/f)\sigma_1\partial_y + \mu_\psi\sigma_2$ and, with the regular tip $b(-L) = 0$, the condition $\theta = 0$ leaves the single zero mode $\chi = (e^{f\mu_\psi y}, 0)$ (DERIVED in the Kohn-Sham record), a chiral four-dimensional fermion localised at the fold for $\mu_\psi > 0$ and at the tip for $\mu_\psi < 0$. The doublets $L, Q$ carry the left-handed zero modes, the singlets $e_R, \nu_R, u_R, d_R$ the right-handed ones (the opposite $\gamma^{x_0}$ component, the zero mode $(0, e^{-f\mu_\psi y})$ with the other fold condition $a(0) = 0$ and the transformed tip condition $a(-L) = 0$: the map $\chi \to \sigma_2\chi$, $\mu \to -\mu$, fold and tip $\theta \to \pi - \theta$ of the Kohn-Sham theorem T3). Gauge-anomaly cancellation of the four-dimensional zero-mode theory is that of one Standard Model generation (STANDARD); anomalies of the eight-dimensional bulk theory are OPEN.
\item The Higgs field is ABSENT. Electroweak symmetry is broken at the tip by a boundary condition: the tip carries a NON-LINEAR sigma field $\Sigma \in SU(2)$, $\Sigma^\dagger\Sigma = 1$, transforming as $\Sigma \to U_L\Sigma U_Y^\dagger$ (INVENTED for this geometry; the structure is that of the electroweak chiral Lagrangian and of Higgsless boundary conditions, STANDARD). $\Sigma$ has no radial mode, hence no Higgs particle; in the unitary gauge $\Sigma = 1$.
\end{itemize}
\[\adjustbox{max width=\linewidth}{$\displaystyle S_{tip} = \int_{y=-L} d^4x\, d^3x_t\sqrt{|g_{tip}|}\Big[\frac{v_\Sigma^2}{4}\mathrm{tr}\big(D_\mu\Sigma^\dagger D^\mu\Sigma\big) - \big(\bar L\,\Sigma\,\Lambda_\ell\,\ell_R + \bar Q\,\Sigma\,\Lambda_q\,q_R + \mathrm{h.c.}\big)\Big]$}\]
with the right-handed doublets $\ell_R = (\nu_R, e_R)$, $q_R = (u_R, d_R)$ and the diagonal coupling matrices $\Lambda_\ell = \mathrm{diag}(\lambda_\nu, \lambda_e)$, $\Lambda_q = \mathrm{diag}(\lambda_u, \lambda_d)$ (the Yukawa structure of the electroweak chiral Lagrangian, gauge invariant because $U(1)_Y$ acts on $\Sigma$ from the right). The first term gives the $W$ and $Z$ their masses with $m_Z = m_W/\cos\theta_W$ at tree level (custodial $SU(2)$ of $\Sigma$); the second gives the fermion masses through the values of the fermion profiles at the tip.
\begin{itemize}
\item QCD (STANDARD): the gluons are the massless $x_0$ zero modes of $G^\alpha_M$; below $\Lambda_{QCD}$ the $u$ and $d$ quarks are confined into protons ($uud$) and neutrons ($udd$), whose masses come from QCD and are inputs of the cosmology (section 7).
\end{itemize}

\subsection*{4.4 The mass mechanism (the author's idea, made precise; DERIVED in \texttt{FIELD\_\allowbreak{}EQUATIONS.\allowbreak{}md} section 4)}
\addcontentsline{toc}{subsection}{4.4 The mass mechanism (the author's idea, made precise; DERIVED in \texttt{FIELD\_\allowbreak{}EQUATIONS.\allowbreak{}md} section 4)}

Every field is expanded on the cell, $\Phi(x) = \sum_n \phi_n(x_1, \dots, x_4)\,u_n(x_0)$, with $u_n$ the eigenfunctions of its $x_0$ operator (periodic with period $\ell_0$, the fold and tip conditions). The four-dimensional mass is the eigenvalue, $m_n^2 = \lambda_n$, and the Rayleigh quotient gives
\[\adjustbox{max width=\linewidth}{$\displaystyle m^2 = \frac{\int dy\, w(y)\,|\partial_y u|^2/f^2 + \dots}{\int dy\, w(y)\,|u|^2},\qquad w(y) = e^{3Hy} \ (\text{the measure } \sqrt{|g|}\,dx_0 \propto e^{3Hy}dy) .$}\]
For a profile localised within a proper length $\Delta s = f\,\Delta y$ the uncertainty relation on the cell gives $\langle \hat p^2\rangle \geq \hbar^2/(4\Delta s^2)$ for the hidden momentum in the orthonormal frame (the author's estimate $m\,c^2 \geq \hbar c/(2\Delta s)$). CORRECTED in version 1.2 (\texttt{ANALYSIS.\allowbreak{}md} sections 3 to 6): in the flat proper measure $v = e^{3Hy/2}u$ the mass is
\[\adjustbox{max width=\linewidth}{$\displaystyle m^2 = \langle\hat p^2\rangle + \Big(\frac{3H}{2f}\Big)^2 - \frac{3H}{f}\,\rho(\text{fold}) + \frac{3H}{f}\,\rho(\text{tips}) + (\text{tip field}) ,$}\]
with $\rho$ the normalised proper density: the warp of the extra times, an attractive point term at the fold and repulsive tip terms (spin 1/2: the bulk mass and the kink of the Jackiw-Rebbi mechanism). The estimate bounds $m$ when the fold term is absent or repulsive: every odd mode and every excitation, $m \geq 3H/(2f)$. For the zero modes the fold term cancels the kinetic energy exactly, and the cancellation is protected by gauge invariance, general covariance and chirality.
\begin{itemize}
\item MASSLESS (no $x_0$ dependence, or an exact zero eigenvalue): the graviton, the photon and the gluons (flat profiles of unbroken gauge generators), the fold-band quanta of $\Psi$ (the fold zero modes, $\varepsilon = e^{-a_4}|k|$).
\item MASSIVE: $W^\pm$ and $Z$ (their profiles bend at the tip, where $\Sigma$ acts), $e$, $\nu$, $u$, $d$ (chiral zero modes localised by $\mu_\psi$, lifted by $\lambda_f$ at the tip: $m_f = \lambda_f\,\zeta_{fL}\zeta_{fR}$ with $\zeta(\mu)^2 = 2f\mu\,e^{-2f\mu L}/(1 - e^{-2f\mu L})$ the normalised value of the zero mode $(e^{f\mu y}, 0)$ at the tip), the bulk-band quanta of $\Psi$ and every excited $x_0$ mode (the Kaluza-Klein tower, $m_n \simeq n\pi/(fL)$ for $n \geq 1$).
\item Because the profiles depend exponentially on $\mu_\psi f L$, order-one parameters produce the observed hierarchy $m_\nu \ll m_e \ll m_u, m_d \ll m_W$ (DERIVED as a fit in \texttt{SOLUTION.\allowbreak{}md}). The hierarchy runs AGAINST the uncertainty estimate: a sharper localisation (larger $\mu_\psi fL$) gives a lighter fermion, and the light masses lie $10^{-3}$ ($W$) to $10^{-16}$ ($\nu$) below $\hbar c/(2\Delta s)$ (\texttt{ANALYSIS.\allowbreak{}md} section 5).
\item Every $x_0$-generated mass depends on the block size $f$: $m \propto 1/f$ for Kaluza-Klein masses and $m \propto \exp(-\mu f L)$ factors for the localised fermions. A deflating block ($h_4 < 0$) therefore makes masses drift in time; the model requires $|h_4| \ll \dot a_4$ after the primordial era (section 7; the bound is computed in \texttt{SOLUTION.\allowbreak{}md}). The extra times keep deflating (never static), but slowly.
\end{itemize}

\subsection*{4.5 The portal interactions of $\Psi$ with the Standard Model (INVENTED; gauge invariant)}
\addcontentsline{toc}{subsection}{4.5 The portal interactions of $\Psi$ with the Standard Model (INVENTED; gauge invariant)}

$\Psi$ is neutral, so every interaction with the Standard Model is through gauge-invariant composite operators. The model contains the lowest ones (eight-dimensional coefficients; the four-dimensional couplings follow from the overlaps of the $x_0$ profiles):

\begin{longtable}{>{\raggedright\arraybackslash}p{0.044\linewidth}>{\raggedright\arraybackslash}p{0.373\linewidth}>{\raggedright\arraybackslash}p{0.171\linewidth}>{\raggedright\arraybackslash}p{0.264\linewidth}}
\toprule
\textbf{} & \textbf{operator} & \textbf{couples $\Psi$ to} & \textbf{cosmological role} \\
\midrule
P1 & $\frac{c_V}{\Lambda_D^{6}}\,(\bar\Psi\gamma^M\Psi)\,(\bar\psi\gamma_M\psi)$, $\psi \in \{L, e_R, Q, u_R, d_R\}$ & $e$, $\nu$, $u$, $d$ (vector currents) & dark matter - baryon and - neutrino scattering \\
P2 & $\frac{c_S}{\Lambda_D^{6}}\,(\bar\Psi\Psi)\,\big(\bar\psi_L\Sigma\psi_R + \mathrm{h.c.}\big)\big|_{tip}$ & $e$, $u$, $d$ (scalar densities) & spin-independent dark matter - nucleon scattering \\
P3 & $\frac{\mu_D}{\Lambda_D^{3}}\,\bar\Psi S^{MN}\Psi\, B_{MN}$ & photon and $Z$ ($B = \cos\theta_W A - \sin\theta_W Z$) & dark dipole: dark matter - photon drag, $Z$ decays \\
P4 & $\frac{c_Z}{\Lambda_D^{4}}\,(\bar\Psi\gamma^M\Psi)\,\frac{i v_\Sigma^2}{2}\mathrm{tr}\big(\Sigma^\dagger D_M\Sigma\,\tau^3\big)\big|_{tip}$ & $Z$ (directly) & $Z$-portal \\
P5 & $\frac{c_D}{\Lambda_D^{7}}\,(\bar\Psi\Psi)\,\big(\bar\psi\, i\gamma^M D_M\psi + \mathrm{h.c.}\big)$ & $W^\pm$, $Z$, $\gamma$, gluons with $e, \nu, u, d$ & charged-current contact vertices \\
\bottomrule
\end{longtable}

\begin{itemize}
\item Charge conservation forbids any gauge-invariant operator in which the neutral $\Psi$ couples to a single $W^\pm$ without a charged pair of Standard-Model fermions; the $W^\pm$ enter the portal through the covariant derivative of P5 (vertices $\bar\Psi\Psi\,W^+\bar e\nu$, $\bar\Psi\Psi\,W^+\bar d u$) and through the gauge self-interactions (DERIVED in \texttt{FIELD\_\allowbreak{}EQUATIONS.\allowbreak{}md} section 5).
\item The coefficients are free parameters bounded by cosmology and by laboratory limits (\texttt{SOLUTION.\allowbreak{}md} section 6). The word CRITICAL of the order is implemented as follows: these are the only couplings through which the dark sector exchanges energy and momentum with the visible sector; without them $\Psi$ interacts with the Standard Model only through gravity.
\end{itemize}

\section*{5. Dimensional reduction (DERIVED at leading order in \texttt{FIELD\_\allowbreak{}EQUATIONS.\allowbreak{}md})}
\addcontentsline{toc}{section}{5. Dimensional reduction (DERIVED at leading order in \texttt{FIELD\_\allowbreak{}EQUATIONS.\allowbreak{}md})}

\begin{itemize}
\item The four-dimensional theory is obtained by integrating over the cell and the extra-time block with the zero-mode profiles. The reduced Planck mass and gauge couplings carry the proper volume $V_\perp = 2\int_{-L}^{0} f\,e^{3Hy}\,dy \times f^3 v_t = 2 f^4 v_t (1 - e^{-3HL})/(3H)$ (2 for the two halves of a cell, $v_t$ the coordinate volume of the extra times; their compactness is ASSUMED, normalisation A of the dark-sector record): $M_{Pl}^2 = V_\perp/\kappa$, $g^2 = \hat g^2/V_\perp$.
\item The extra times are TIME-LIKE. Modes with extra-time momentum $q$ have $\omega^2 = M^2 + k^2 - q^2$ and are tachyonic for $q > M$ (DERIVED in the dark-sector record); the model keeps only the good sector $q = 0$ (ASSUMED), as the Kohn-Sham record does. The stability of the full theory is OPEN.
\end{itemize}

\section*{6. The cosmological inventory}
\addcontentsline{toc}{section}{6. The cosmological inventory}

\begin{longtable}{>{\raggedright\arraybackslash}p{0.221\linewidth}>{\raggedright\arraybackslash}p{0.179\linewidth}>{\raggedright\arraybackslash}p{0.170\linewidth}>{\raggedright\arraybackslash}p{0.281\linewidth}}
\toprule
\textbf{component} & \textbf{field} & \textbf{$x_0$ mode} & \textbf{equation of state} \\
\midrule
photons (CMB) & $A_\mu$ zero mode & flat, massless & $w = 1/3$ \\
neutrinos (CNB) & $\nu_L$ zero mode & localised, $m_\nu \lesssim 0.1$ eV & relativistic, then non-relativistic \\
baryons & $u, d \to p, n$ & localised, QCD & $w = 0$ \\
electrons & $e$ & localised & $w = 0$ after $e^\pm$ annihilation \\
dark matter (CONDITIONAL, \texttt{ANALYSIS.\allowbreak{}md} section 8) & $\Psi$ bulk band, lightest odd level & massive, $\varepsilon \geq m$ & $w = 0$ on the static block; $-0.073$ while the block deflates \\
dark radiation & $\Psi$ fold band & massless zero modes & $w = 1/3$ ($\Delta N_{eff}$) \\
dark condensate & $\bar\Psi\Psi$ & homogeneous & $w_c = \lambda S/(2m + \lambda S)$, drifting; dust-like on the static block, NOT dark energy \\
hidden curvature & geometry, $6H^2 k$, $k = e^{2a_4}/b_4$ & - & acts like a cosmological term when $k$ is constant \\
dark energy & $\delta\Lambda$ (geometric, tuned) & - & $w = -1$; no $\Psi$ component qualifies (\texttt{ANALYSIS.\allowbreak{}md} section 9) \\
\bottomrule
\end{longtable}

\section*{7. What the model predicts and what it requires (statements to be computed)}
\addcontentsline{toc}{section}{7. What the model predicts and what it requires (statements to be computed)}

\begin{itemize}
\item The background equations (DERIVED from the record, \texttt{FIELD\_\allowbreak{}EQUATIONS.\allowbreak{}md} section 2): in Einstein gravity the $x_4$ constraint is $3\dot a_4^2 + 12\,\dot a_4 h_4 + 6h_4^2 - 6H^2 k = \kappa\rho + \Lambda$ with $k = e^{2a_4}/b_4$: a Friedmann equation with the extra term $12\dot a_4 h_4 + 6h_4^2$ of the deflating block and the hidden-curvature term $6H^2k$, which is constant when the block is static and acts like a cosmological constant.
\item The masses: the $W$, $Z$ and fermion masses are fitted with order-one parameters; the photon, gluon and graviton are exactly massless (\texttt{SOLUTION.\allowbreak{}md} section 4).
\item CMB, CNB and BAO: computed in the four-dimensional effective theory at order $\epsilon$ (\texttt{SOLUTION.\allowbreak{}md} section 5): the neutrino temperature $T_\nu = (4/11)^{1/3}T_\gamma$ (STANDARD, unchanged as long as $\Psi$ does not heat the neutrinos), the sound horizon at the drag epoch, the angular acoustic scale and the BAO scale, and the bounds that the deflation rate, the dark radiation and the portal couplings must respect.
\end{itemize}

\section*{8. Assumptions and open problems (honesty rule)}
\addcontentsline{toc}{section}{8. Assumptions and open problems (honesty rule)}

\begin{itemize}
\item ASSUMED: the regular tip at $y = -L$; the compact extra times (normalisation A) and the good sector $q = 0$; the kink bulk masses; the tip sigma field; the portal operators; one generation; the zero-mode truncation at order 1; (1.1) the block stabiliser and the exact fold-reflection parity that keeps the odd-parity dark matter stable ((1.2) exact to $2\times10^{-27}$ in amplitude for a lifetime above $10^{26}$ s, \texttt{ANALYSIS.\allowbreak{}md} section 8.3).
\item REQUIRED (1.2, \texttt{ANALYSIS.\allowbreak{}md}): the exponential deflation of the primordial era ends long before nucleosynthesis (the electron mass changes like $f^{-22}$); the portal coefficients of the $\Psi$ fold band and of the $\Psi$ dark matter differ beyond their overlaps, or the one-half count $g_F = 4$ holds (otherwise the $N_{eff}$ portal makes the dark matter under-abundant and puts it at 1.6 times the LZ limit).
\item CHANGED (1.1): $w_3 = 1/768$ and $\Lambda = -24H^2 + \delta\Lambda$ (section 4.1). The solution also finds what the model requires of the data: $m_{KK} \approx 121$ TeV from $\rho_0$, and, with one generation, a decoupling of the $\Psi$ fold band inside the QCD crossover for $N_{eff} \approx 3$ (\texttt{SOLUTION.\allowbreak{}md} sections 4 and 6).
\item OPEN: the stability of a theory with three time-like extra dimensions (ghosts and tachyons beyond the good sector); perturbative unitarity of $WW$ scattering without a Higgs (it holds up to the Kaluza-Klein scale $\sim \pi/(fL)$ only if that scale lies below about $4\pi v_\Sigma \approx 3$ TeV, as in all Higgsless models); a dynamical mechanism that slows the deflation of the block after the primordial era; the origin of the second and third generations; the derivation of the portal coefficients; the non-adiabatic evolution of $\Psi$; (1.2) the annihilation that fixes the dark-matter abundance ($6.5\times10^{-26}$ cm$^3$/s, 37 \% of the s-wave unitarity limit).
\item Consistency with QCD and the electroweak theory is at the level of the gauge group, the representations, anomaly cancellation and the electroweak chiral Lagrangian (no Higgs); precision electroweak constraints on Higgsless models (the $S$ parameter) apply and are OPEN for this geometry.
\end{itemize}

\section*{9. Analysis of the mass mechanism and of the dark associations (version 1.2; DERIVED in \texttt{ANALYSIS.\allowbreak{}md})}
\addcontentsline{toc}{section}{9. Analysis of the mass mechanism and of the dark associations (version 1.2; DERIVED in \texttt{ANALYSIS.\allowbreak{}md})}

\begin{itemize}
\item The mass mechanism. Every four-dimensional mass is an eigenvalue of an $x_0$ operator of a wave function periodic with $\pi/(6H)$: the author's requirement holds. Without a cutoff each period is a whole line in proper length, and every function of $y$ is periodic; the content of the requirement lies in the parities and boundary terms at the fold and the tips. The operator factorises as $A^\dagger A$: the mass is the $x_0$ dependence of the profile RELATIVE to the geometric zero mode. A mode localised as the geometry dictates is massless, and every deviation costs mass.
\item The uncertainty relation is a true lower bound for the $x_0$ excitations ($m \geq 3H/(2f) = 99.4$ TeV; the estimate is exceeded by factors 2.6 to 7.5). The dark matter is one of them: 78.9 \% of its $\varepsilon^2$ is $x_0$ momentum. The Kohn-Sham stress tensor confirms this: its hidden pressure $p_0/\rho = 0.6575394$ equals $-d\ln M_{DM}/d\ln f$ to $10^{-9}$.
\item The uncertainty relation is not the origin of the light masses. The photon, gluon and graviton zero mode is localised within $1.4\times10^{-21}$ m and massless. The $W$ mass is 99.87 \% tip term. The fermion masses fall as the localisation sharpens. Taken literally as a principle, the Heisenberg argument would put every massive particle above 99 TeV.
\item Constancy: every $x_0$-generated mass tracks $f$ ($m_e \propto f^{-22}$, $\alpha, G_N \propto f^{-4}$); the late block deflation $h_4/A \approx -10^{-93}$ is far within the atomic-clock bound $1.2\times10^{-8}$, and the primordial deflation must end before nucleosynthesis.
\item Dark matter: CONDITIONAL. The odd bulk quanta are cold on the static block, dark, collisionless and massive. They need three things: an exact fold parity (ASSUMED), an annihilation of natural cell strength ($G \approx 1.4/H^2$, OPEN), and portal coefficients that are not universal across $\Psi$ modes. The record's computed states are dark radiation.
\item Dark energy: NOT SUPPORTED. No positive-energy $\Psi$ state has $w < 0$ on the static block. The $\lambda < 0$ condensate tuned to $w_0 = -1$ has $w_a = -6$, and its energy turns negative at $z = 0.26$. The vacuum energy is the cosmological-constant problem. The dark energy of UM-1 is $\delta\Lambda$.
\end{itemize}

\end{document}
